\subsection{Exact targets in synthetic transaction histories}\label{sec:batch}

The file \texttt{HI-Small\_Trans.csv} from version~8 of IBM's synthetic
anti-money-laundering dataset~\citep{altman2023} contains 5,078,345 transactions.
The sampling unit was a destination
account and transaction minute. An earlier audit fixed 1,024 queries by
seeded hash order from the 3,854,758 distinct units with a full 24-hour
lookback, without using labels or structural eligibility. These queries
contain 1,023 distinct roots. Each context includes every transaction in
$[t-1440,t)$ between accounts within incoming static distance three of the
root in that window. The extraction retained all such events, including loops, tied
times and parallel contacts; no conflicting edge was removed to make a
context eligible. The cohort had been studied previously, and queries may
share events. The results describe these queries rather than independent
trials or an unseen test set.

\paragraph{Coverage.}
All 1,024 contexts were extracted and their baseline counts determined.
The structural check was applied to the original induced context.
Table~\ref{tab:batch-coverage} retains empty, unsupported, and
resource-limited queries in the denominator. Of 232 queries with
baseline $W_3>0$, 77 fit the model and 68 received complete profiles.
The nine skipped eligible positive queries had a median of 91 events
compared with 12 for the 68 profiled queries.

\begin{center}
\begin{tabular}{lrrr}
\toprule
Query disposition & $W_3=0$ & $W_3>0$ & Total\\
\midrule
Empty context & 179 & 0 & 179\\
Outside the required tree structure & 125 & 155 & 280\\
Eligible, above the 32-event cap & 34 & 9 & 43\\
Eligible, all requested profiles completed & 454 & 68 & 522\\
\midrule
Total & 792 & 232 & 1,024\\
\bottomrule
\end{tabular}
\refstepcounter{table}\label{tab:batch-coverage}\\[3pt]
{\small Table \thetable. Disposition of all sampled account/minute queries.}
\end{center}

\paragraph{Two questions asked of the same profile.}
Every event was optional with unit deletion cost. Given baseline counts
$N_2,N_3$, set $T=\lceil pN_2/100\rceil$ for $p\in\{100,95,90\}$ and
$U=\lfloor 3N_3/4\rfloor$. For each $T$, the evaluation compared the goals $W_3\le U$ and
$W_3=U$, both subject to $W_2\ge T$. Budgets 0, 1 and 3 were applied to
the minimum deletion cost. The 454 profiled queries with $N_3=0$ satisfy
both goals without intervention; the comparison below therefore concerns
the 68 profiled queries with $N_3>0$.

\begin{center}
\begin{tabular}{lrrr}
\toprule
 & \multicolumn{2}{c}{Attainable within three deletions}
 & Unattainable at any cost\\
\cmidrule(lr){2-3}\cmidrule(l){4-4}
Preservation & $W_3\le U$ & $W_3=U$ & $W_3=U$\\
\midrule
$100\%$ & 68 & 13 & 55\\
$95\%$  & 68 & 14 & 54\\
$90\%$  & 68 & 15 & 53\\
\bottomrule
\end{tabular}
\refstepcounter{table}\label{tab:batch-goals}\\[3pt]
{\small Table \thetable. Each row refers to the same 68 positive, fully profiled queries.}
\end{center}

At full preservation, the least cost for the upper bound was one deletion in
59 queries, two in eight, and three in one. In 55 of the 68, no deletion set
could attain the exact same target while preserving every pair, even without
a deletion budget. At $R=0$, every middle or final event belonging to a
baseline pair must survive. Deleting a first event then removes all triples
extending it as one group; different first events remove disjoint groups.
Attainable triple losses are sums of these group sizes, so some intermediate
integer losses can be absent. For example, in the preidentified pilot
query \texttt{q0011}, preserving seven pairs permits either two triples at
cost zero or zero triples at cost one. Exactly one triple is unattainable.
Rounding also matters: on 50 of the 68 positive queries all three
percentages require the same integer $T$. The rows of
Table~\ref{tab:batch-goals} are paired comparisons, not 204 independent
observations.

\paragraph{Verification and scope.}
Deduplicating equal values of $T$ gave 653 planned tables among the 565
eligible nonempty queries. The 32-event cap excluded 90 tables before
optimization; these have unknown outcomes. There were no solver or oracle
timeouts. The Mac run reported witness replay for all 935 returned entries
across the 563 completed tables. For 528 of those tables (at most 16
optional events), exhaustive enumeration agreed on every attainable pair
and its minimum cost. The other 35 had witness replay but no exhaustive
comparison. The run provides evidence about local count sensitivity and
model coverage on synthetic transactions. It does not measure classification
accuracy, changes in fraud decisions, or comparative runtime.
